\chapter{Cornea Transplant}

\section*{What Is This Surgery?}
Cornea transplant, also known as keratoplasty, is a sophisticated surgical procedure that involves replacing a diseased, scarred, or damaged cornea with healthy donor corneal tissue. The cornea, the transparent, dome-shaped outermost layer at the front of the eye, is critically responsible for focusing light onto the retina and protecting the intraocular structures. This intervention is primarily performed to restore vision that has been compromised by corneal opacification, distortion, or thinning, or to alleviate severe ocular pain caused by intractable corneal disease. It represents a vital therapeutic option for a spectrum of conditions that are not amenable to correction through less invasive medical or surgical means, thereby significantly enhancing the quality of life for patients afflicted with corneal blindness or profound discomfort. The enduring success of corneal transplantation is underpinned by the cornea's unique immunological properties and continuous advancements in microsurgical techniques and comprehensive postoperative care protocols.

\section*{Alternative Names}
\begin{itemize}
  \item Keratoplasty
  \item Corneal Graft
  \item Full-Thickness Corneal Transplant
  \item Partial-Thickness Corneal Transplant
  \item Endothelial Keratoplasty
\end{itemize}

\section*{Relevant Surgical Anatomy}
The cornea is a remarkable, avascular, transparent tissue forming the anterior one-sixth of the fibrous tunic of the eye. Its average diameter is approximately 11.5 mm, with a central thickness of about 550 micrometers, gradually increasing towards the periphery. Structurally, the cornea is composed of six distinct layers, from anterior to posterior: the corneal epithelium, Bowman's layer, corneal stroma, Descemet's membrane, and the corneal endothelium. A recently identified sixth layer, Dua's layer, lies between the posterior stroma and Descemet's membrane. The epithelium, a stratified squamous non-keratinized layer, provides a smooth optical surface and a barrier against infection. Bowman's layer is a strong, acellular collagenous layer. The stroma, comprising about 90\% of corneal thickness, consists of highly organized collagen lamellae and keratocytes, contributing significantly to corneal strength and transparency. Descemet's membrane is the basement membrane of the endothelium, and the endothelium itself is a single layer of hexagonal cells responsible for actively pumping fluid out of the stroma, maintaining corneal deturgescence and transparency.

The cornea is entirely avascular, receiving its nutrients from the aqueous humor, the limbal capillaries (at the corneoscleral junction), and tear film. This avascularity is a key factor in its relative immune privilege. Sensory innervation is exceptionally rich, primarily supplied by the long ciliary nerves, branches of the ophthalmic division of the trigeminal nerve (CN V1). These nerves enter the cornea radially from the limbus, forming a dense subepithelial plexus, making the cornea one of the most sensitive tissues in the body. Lymphatic drainage is virtually absent within the cornea itself, further contributing to its immune privilege, though lymphatic vessels are present at the limbus. Surgical landmarks include the limbus, which is the transitional zone between the cornea and sclera, and the anterior chamber, the space between the cornea and iris/lens, filled with aqueous humor. Understanding these anatomical relationships is paramount for precise surgical planning and execution, particularly in selecting the appropriate keratoplasty technique to target specific diseased layers while preserving healthy ones.

\section{Surgical Principle \& Rationale}
The fundamental surgical principle of corneal transplantation is the meticulous replacement of diseased or dysfunctional corneal tissue with healthy, optically clear donor tissue to restore the cornea's primary functions: optical clarity and structural integrity. The rationale for this procedure stems from the cornea's unique biological characteristics, notably its avascularity and the immune-privileged environment of the anterior chamber (anterior chamber associated immune deviation - ACAID), which collectively contribute to a relatively lower risk of allograft rejection compared to other solid organ transplants.

Modern keratoplasty techniques are precisely tailored to the specific corneal layer affected by disease. For full-thickness pathologies such as severe scarring, advanced keratoconus, or corneal perforation, Penetrating Keratoplasty (PKP) is employed, replacing all corneal layers to restore both optical and structural integrity. When only the anterior stromal layers are diseased, as in certain forms of keratoconus or superficial scars, Deep Anterior Lamellar Keratoplasty (DALK) is preferred, preserving the patient's healthy Descemet's membrane and endothelium. This lamellar approach significantly reduces the risk of endothelial graft rejection, which is the most common cause of graft failure. Conversely, for conditions primarily affecting the corneal endothelium, such as Fuchs' endothelial dystrophy or pseudophakic bullous keratopathy, endothelial keratoplasty techniques like Descemet's Stripping Endothelial Keratoplasty (DSEK) and Descemet's Membrane Endothelial Keratoplasty (DMEK) are utilized. These procedures selectively replace only the diseased endothelial layer and Descemet's membrane, preserving the patient's healthy anterior cornea. The rationale for these selective techniques is to minimize surgical invasiveness, accelerate visual recovery, and further reduce the risk of immune rejection by transplanting less tissue, particularly avoiding the highly immunogenic stromal keratocytes.

\section{History \& Surgical Evolution}
The ambitious concept of corneal transplantation traces its origins to the early 19th century, with pioneering, albeit largely unsuccessful, experimental attempts by surgeons such as Franz Reisinger in 1824. These initial endeavors were hampered by rudimentary surgical techniques and an incomplete understanding of tissue rejection. The pivotal moment in the history of corneal surgery arrived in 1905 when Austrian ophthalmologist Eduard Zirm performed the first truly successful full-thickness corneal transplant (Penetrating Keratoplasty). Using tissue from an enucleated eye, Zirm successfully restored vision in a patient suffering from severe corneal burns, marking the dawn of modern corneal transplantation.

The mid-20th century witnessed significant refinements in PKP techniques, largely driven by pioneers like Ramón Castroviejo, who developed specialized trephines and advanced suturing methods that dramatically improved graft survival and visual outcomes. The introduction of corticosteroids in the 1950s proved revolutionary, providing a critical pharmacological tool for managing postoperative inflammation and mitigating graft rejection. The late 20th and early 21st centuries ushered in a paradigm shift with the development of lamellar and endothelial keratoplasty techniques. Gerrit Melles introduced posterior lamellar keratoplasty in the late 1990s, which subsequently evolved into Descemet's Stripping Endothelial Keratoplasty (DSEK) by Mark Gorovoy and Francis Price Jr. Further refinement by Melles and others led to Descemet's Membrane Endothelial Keratoplasty (DMEK), a highly selective procedure involving the transplantation of only Descemet's membrane and its attached endothelium. These advancements have progressively moved corneal surgery towards more selective, less invasive, and anatomically precise tissue replacement, optimizing outcomes for specific corneal pathologies.

\section{Clinical Indications}
Cornea transplant is indicated for a diverse range of corneal pathologies that cause significant visual impairment or intractable ocular pain and are not amenable to medical management or less invasive surgical interventions.

\begin{itemize}
  \item Advanced keratoconus or pellucid marginal degeneration that has progressed to a stage where vision cannot be adequately corrected with rigid gas permeable contact lenses or other interventions like corneal collagen cross-linking.
  \item Fuchs' endothelial dystrophy, a hereditary condition leading to progressive endothelial cell loss, resulting in chronic corneal edema and bullous keratopathy.
  \item Corneal scarring caused by severe trauma, infectious keratitis (e.g., herpes simplex keratitis, bacterial or fungal infections), chemical or thermal burns, or complications from previous ocular surgeries.
  \item Pseudophakic or aphakic bullous keratopathy, characterized by endothelial decompensation and corneal edema following cataract extraction.
  \item Corneal perforation or severe corneal thinning (descemetocele) due to infection, autoimmune diseases (e.g., rheumatoid arthritis-associated keratitis), or traumatic injury.
  \item Failure of a previous corneal graft, necessitating re-transplantation due to rejection, primary graft failure, or other complications.
  \item Severe corneal edema from other etiologies, such as congenital hereditary endothelial dystrophy in infants or children.
  \item Painful bullous keratopathy in eyes with poor visual potential, where the primary goal of surgery is to alleviate chronic pain and discomfort.
  \item Corneal dystrophies (e.g., lattice dystrophy, granular dystrophy) that cause significant visual compromise due to stromal opacification.
\end{itemize}

\section{Clinical Positioning}
Cornea transplant can serve as both a primary, first-line definitive treatment and a secondary or salvage procedure, depending on the specific underlying corneal pathology, its severity, and the patient's overall ocular health. For conditions such as advanced Fuchs' endothelial dystrophy, severe, full-thickness corneal scarring, or acute corneal perforation, keratoplasty is often considered a primary intervention. In these scenarios, vision loss is typically significant and irreversible through medical management or less invasive procedures, making the transplant the most direct and effective means to restore optical clarity and structural integrity.

Conversely, for other conditions, such as keratoconus, corneal transplantation is frequently a secondary or tertiary procedure. Initial management for keratoconus typically involves non-surgical options like rigid gas permeable contact lenses to correct irregular astigmatism, or corneal collagen cross-linking to halt disease progression. A transplant is reserved for cases where these primary treatments fail, the disease progresses to a stage where contact lens wear is no longer tolerable or effective, or severe scarring develops. Similarly, a corneal transplant may be a secondary or salvage procedure for complications arising from previous ocular surgeries, such as pseudophakic bullous keratopathy, or for the failure of an earlier corneal graft due to rejection or other issues. In these contexts, the transplant aims to salvage vision or ocular integrity after initial treatments have proven insufficient or failed.

\section{Surgical Technique \& Procedure}
Cornea transplant surgery is a highly specialized microsurgical procedure typically performed under local anesthesia with intravenous sedation, though general anesthesia may be preferred for anxious patients, pediatric cases, or complex re-transplantations. The patient is positioned supine, and the head is carefully stabilized. The specific operative steps vary significantly based on the type of keratoplasty performed.

\paragraph{}
**1. Penetrating Keratoplasty (PKP) - Full-Thickness Transplant:**
\begin{enumerate}
  \item  **Anesthesia and Preparation:** Local anesthetic injection (e.g., peribulbar block) with sedation, or general anesthesia. The eye is prepped and draped, and a lid speculum is inserted.
  2.  **Recipient Cornea Trephination:** A circular trephine of appropriate diameter (typically 7.5-8.5 mm) is used to create a full-thickness incision in the diseased recipient cornea, removing the central corneal button.
  3.  **Donor Cornea Preparation:** A healthy donor cornea, provided by an eye bank and precisely matched in size (often 0.25-0.5 mm larger than the recipient opening to ensure a tight fit and reduce postoperative hyperopia), is prepared. The donor button is punched from the endothelial side.
  4.  **Donor Graft Placement:** The donor corneal button is carefully placed into the recipient bed.
  5.  **Suturing:** The donor graft is secured with numerous fine sutures. This typically involves four cardinal interrupted sutures initially, followed by 12-16 additional interrupted sutures, or a continuous running suture, or a combination. Suture tension is meticulously adjusted to minimize postoperative astigmatism.
  6.  **Anterior Chamber Reformation:** The anterior chamber is reformed with balanced salt solution, and intraocular pressure is checked.
  7.  **Closure:** The conjunctiva is not typically closed over the graft. A topical antibiotic and steroid are applied, and an eye shield is placed.
\end{enumerate}

\paragraph{}
**2. Deep Anterior Lamellar Keratoplasty (DALK) - Anterior Stromal Transplant:**
\begin{enumerate}
  \item  **Anesthesia and Preparation:** Similar to PKP.
  2.  **Partial-Thickness Trephination:** A partial-thickness trephination is performed on the recipient cornea, extending to approximately 80-90\% of stromal depth, leaving Descemet's membrane intact.
  3.  **Stromal Dissection (Big Bubble Technique):** Air or viscoelastic is injected into the deep stroma to create a "big bubble" separating Descemet's membrane from the overlying stroma, facilitating a clean dissection.
  4.  **Removal of Diseased Stroma:** The diseased anterior stromal layers are carefully dissected and removed, exposing the bare Descemet's membrane.
  5.  **Donor Cornea Preparation:** A donor corneal cap, with its endothelium and Descemet's membrane meticulously removed, is prepared.
  6.  **Donor Graft Placement and Suturing:** The donor stromal cap is placed onto the exposed recipient Descemet's membrane and secured with interrupted or running sutures, similar to PKP.
  7.  **Closure:** Topical medications and an eye shield are applied.
\end{enumerate}

\paragraph{}
**3. Endothelial Keratoplasty (DSEK/DMEK) - Endothelial Transplant:**
\begin{enumerate}
  \item  **Anesthesia and Preparation:** Similar to PKP, often with a retrobulbar block for akinesia.
  2.  **Small Incision:** A small limbal incision (2-3 mm) is created, along with one or two paracentesis incisions.
  3.  **Descemet's Membrane Stripping:** The diseased Descemet's membrane and endothelium are carefully stripped from the posterior surface of the recipient cornea using a reverse Sinskey hook or similar instrument.
  4.  **Donor Tissue Preparation:** For DSEK, a donor corneal lamella (stroma, Descemet's, endothelium) is prepared. For DMEK, a very thin graft consisting only of Descemet's membrane and endothelium is meticulously isolated.
  5.  **Donor Tissue Insertion:** The prepared donor tissue is folded or rolled and inserted into the anterior chamber through the small incision.
  6.  **Donor Tissue Unfolding and Positioning:** The donor tissue is carefully unfolded and positioned against the posterior recipient stroma.
  7.  **Air/Gas Bubble Injection:** An air or gas bubble (e.g., SF6) is injected into the anterior chamber to press the donor tissue against the recipient stroma, facilitating adhesion.
  8.  **Closure:** The small incision is self-sealing or closed with a single suture. Topical medications and an eye shield are applied. The air/gas bubble acts as a temporary internal splint, requiring specific postoperative head positioning.
\end{enumerate}

\section*{Preoperative Workup \& Preparation}
Thorough preoperative workup and preparation are paramount to optimize the success of a corneal transplant and minimize potential complications. This comprehensive evaluation ensures the patient is medically fit for surgery and that the ocular surface is in the best possible condition. The process typically begins with a detailed ophthalmic examination, including visual acuity, intraocular pressure measurement, slit-lamp biomicroscopy, corneal topography and pachymetry to assess corneal shape and thickness, specular microscopy to evaluate endothelial cell density, and anterior segment optical coherence tomography (OCT) for precise layer assessment.

\paragraph{}
In addition to the ocular assessment, the following preparations are crucial:
\begin{itemize}
  \item **Systemic Medical Evaluation:** A comprehensive medical history and physical examination are performed. Routine blood tests (e.g., complete blood count, electrolytes, coagulation profile), electrocardiogram (EKG), and chest X-ray may be ordered, particularly for patients undergoing general anesthesia or those with significant systemic comorbidities, to assess surgical fitness.
  \item **Infection Screening and Treatment:** Any active ocular or systemic infections (e.g., blepharitis, conjunctivitis, dacryocystitis, or systemic viral infections) must be identified and rigorously treated prior to surgery to mitigate the risk of postoperative endophthalmitis or graft infection.
  \item **Medication Review and Adjustment:** All current medications are reviewed. Anticoagulants or antiplatelet agents (e.g., warfarin, aspirin, clopidogrel) may need to be temporarily discontinued or adjusted under the guidance of the patient's primary care physician or cardiologist to minimize bleeding risk.
  \item **NPO Status:** Patients are typically instructed to be nil per os (NPO) for 6-8 hours before surgery to prevent aspiration during anesthesia.
  \item **Prophylactic Antibiotics:** Topical broad-spectrum antibiotics are often prescribed for several days leading up to the surgery to reduce the bacterial load on the ocular surface.
  \item **Informed Consent:** A detailed discussion with the patient and their family regarding the indications, risks, benefits, alternatives, and expected postoperative course of the procedure is conducted, and comprehensive informed consent is obtained.
  \item **Donor Tissue Procurement and Screening:** The eye bank meticulously screens and prepares donor corneal tissue, adhering to strict quality and safety standards, including serological testing for infectious diseases (e.g., HIV, hepatitis).
  \item **Ocular Surface Optimization:** Pre-existing ocular surface diseases, such as severe dry eye syndrome or chronic blepharitis, are aggressively treated to promote optimal epithelial healing postoperatively.
  \item **Patient Education:** Patients undergoing endothelial keratoplasty are specifically educated on the critical importance of postoperative head positioning to ensure proper graft adhesion.
\end{itemize}

\section*{Contraindications \& Precautions}
Careful consideration of contraindications and precautions is essential to ensure patient safety, optimize the likelihood of graft survival, and achieve meaningful visual improvement.

\paragraph{}
**Absolute Contraindications:**
\begin{itemize}
  \item Active, uncontrolled ocular infection (bacterial, viral, fungal, or protozoal keratitis or endophthalmitis) that has not been adequately treated and resolved.
  \item Uncontrolled, severe glaucoma with significant optic nerve damage and poor prognosis, as the surgery and subsequent steroid use may exacerbate intraocular pressure.
  \item Severe, unmanageable ocular surface disease (e.g., severe dry eye syndrome, aniridia-associated keratopathy, extensive chemical burns with persistent epithelial defects) that precludes proper epithelial healing and graft survival.
  \item Significant systemic illness or unstable medical condition (e.g., recent myocardial infarction, uncontrolled diabetes, severe respiratory compromise) that renders the patient unfit for safe administration of anesthesia or recovery from surgery.
  \item Lack of potential for visual improvement due to severe underlying retinal or optic nerve disease (e.g., advanced macular degeneration, optic atrophy), making the corneal transplant futile.
\end{itemize}

\paragraph{}
**Relative Contraindications and Precautions:**
\begin{itemize}
  \item Advanced corneal neovascularization, which significantly increases the risk of graft rejection by providing immune cells direct access to the donor tissue.
  \item History of multiple previous graft failures, indicating a highly sensitized immunological environment and a poorer prognosis for subsequent grafts.
  \item Uncontrolled chronic ocular inflammation (e.g., chronic uveitis, ocular cicatricial pemphigoid), which can lead to persistent inflammation and graft failure.
  \item Poor patient compliance with the rigorous postoperative medication regimen (especially topical corticosteroids) and demanding follow-up schedule, which is critical for graft survival.
  \item Monocular patient (only one functional eye), where the inherent risks of surgery to the only seeing eye must be meticulously weighed against the potential benefits.
  \item Pre-existing glaucoma, which necessitates careful monitoring and aggressive management of intraocular pressure postoperatively, particularly with prolonged steroid use.
  \item History of herpes simplex keratitis, requiring prophylactic oral antiviral medication postoperatively to prevent viral reactivation and graft infection.
  \item Significant psychological or cognitive impairment that may hinder adherence to postoperative instructions.
\end{itemize}

\section*{Complications \& Risks}
While corneal transplantation is generally a safe and highly effective procedure, like all surgical interventions, it carries a spectrum of potential risks and complications. These can be broadly categorized into general surgical risks and procedure-specific risks.

\paragraph{}
**General Surgical Complications:**
\begin{itemize}
  \item **Bleeding:** Intraoperative or postoperative hemorrhage, such as hyphema (blood in the anterior chamber), which can obscure vision and increase intraocular pressure.
  \item **Infection:** Ocular infection, particularly endophthalmitis (infection inside the eye), is a rare but devastating complication that can lead to severe vision loss or even loss of the eye.
  \item **Anesthesia Complications:** Risks associated with local or general anesthesia, including nausea, vomiting, allergic reactions, respiratory depression, or cardiovascular events.
\item **Pain:** Postoperative pain, though usually manageable with oral analgesics.
\end{itemize}

\paragraph{}
**Procedure-Specific Risks:**
\begin{itemize}
  \item **Graft Rejection:** The most common and significant complication, where the recipient's immune system recognizes the donor tissue as foreign and mounts an immune response. It can be acute or chronic and, if not promptly treated, can lead to irreversible graft failure.
  \item **Primary Graft Failure:** The donor cornea fails to clear or function properly from the outset, often due to poor donor tissue quality, prolonged preservation time, or significant surgical trauma.
  \item **High Astigmatism:** Especially prevalent after Penetrating Keratoplasty (PKP), where irregular healing or uneven suture tension can cause significant refractive error, necessitating glasses, contact lenses, or further surgical intervention (e.g., relaxing incisions, excimer laser).
  \item **Glaucoma:** Elevated intraocular pressure (IOP) can occur due to inflammation, steroid-induced ocular hypertension, or changes in anterior chamber anatomy, potentially leading to optic nerve damage and permanent vision loss.
  \item **Retinal Detachment:** A rare but serious complication that can occur due to surgical trauma, changes in intraocular fluid dynamics, or vitreous prolapse.
  \item **Wound Leak or Dehiscence:** Inadequate wound closure, particularly after PKP or DALK, can lead to leakage of aqueous humor, increasing the risk of infection, hypotony, and anterior chamber collapse.
  \item **Cataract Formation or Progression:** The development or worsening of cataracts (opacification of the natural lens) can occur, particularly in older patients or as a side effect of prolonged topical steroid use.
  \item **Persistent Epithelial Defects:** Failure of the recipient epithelium to heal over the donor graft, leading to chronic discomfort, increased risk of infection, and stromal melting.
  \item **Suture-Related Complications:** In PKP/DALK, loose, broken, or infected sutures can cause discomfort, induce astigmatism, or serve as a nidus for infection.
  \item **Dislocation of Endothelial Graft:** In DSEK/DMEK, the donor endothelial tissue may detach from the recipient cornea, requiring a "re-bubble" procedure to re-adhere the graft.
  \item **Endothelial Cell Loss:** A natural process accelerated by surgery, which can lead to late graft failure if the cell count drops below a critical threshold.
\end{itemize}

\section*{Postoperative Recovery Timeline}
Immediately following corneal transplant surgery, the patient is transferred to a Post-Anesthesia Care Unit (PACU) for close monitoring as they emerge from anesthesia. An eye shield is typically placed over the operated eye to provide protection against accidental trauma. Pain management and anti-emetic medications are administered as needed. Most patients are discharged home the same day or the following morning, depending on the type of anesthesia and their recovery status.

\paragraph{}
During the first 24-48 hours, strict adherence to postoperative instructions is crucial. For patients undergoing endothelial keratoplasty (DSEK/DMEK), specific head positioning (e.g., lying supine for 75-90\% of the day) is often required to ensure the intraocular air or gas bubble presses the donor tissue against the recipient cornea, facilitating proper adhesion. Topical medications, including broad-spectrum antibiotics to prevent infection, potent corticosteroids to suppress inflammation and prevent rejection, and cycloplegics to reduce pain and stabilize the iris, are prescribed and must be administered diligently. Patients are advised to avoid rubbing the eye, heavy lifting, bending over, or any activities that could significantly increase intraocular pressure. Frequent follow-up visits, often daily or weekly initially, are scheduled to monitor graft clarity, intraocular pressure, and detect any early signs of complications. Over the first 1-2 weeks, the steroid dosage may be gradually tapered based on the eye's inflammatory response. Long-term, topical corticosteroids are typically continued for several months to years, especially for PKP, to maintain immune suppression and prevent rejection. For PKP and DALK, sutures remain in place for months to years and may be selectively removed to reduce astigmatism. Visual recovery is gradual, often taking several weeks to months for endothelial grafts and considerably longer (up to a year or more) for full-thickness grafts due to the healing process and astigmatism stabilization.

\section*{Surgical Findings \& Pathology}
During a corneal transplant, the surgeon meticulously documents the intraoperative findings, which include the precise dimensions of the trephination, the clarity and integrity of the host cornea's remaining layers (in lamellar procedures), the presence of any synechiae, iris abnormalities, or vitreous prolapse. The condition of the diseased cornea, such as the extent of scarring, edema, vascularization, or thinning, is also noted. For endothelial keratoplasty, the ease of stripping Descemet's membrane and the successful unfolding and positioning of the donor graft are critical findings.

\paragraph{}
If a full-thickness or anterior lamellar section of the recipient cornea is excised, it is sent for histopathological examination. The pathology report serves to confirm the original clinical diagnosis, providing definitive microscopic evidence of conditions such as keratoconus (e.g., thinning of the stroma, breaks in Bowman's layer), Fuchs' endothelial dystrophy (e.g., guttae, endothelial cell loss), corneal scarring (e.g., disorganized collagen, inflammatory cells), or infectious keratitis (e.g., presence of microorganisms, inflammatory infiltrate). The report also assesses the quality and integrity of the donor tissue, ensuring it meets the required standards for transplantation. This pathological confirmation is vital for understanding the disease process and guiding long-term management.

\section*{Expected Postoperative Course}
A normal and successful postoperative course following a corneal transplant is characterized by a gradual and progressive improvement in visual acuity, coupled with a stable and clear corneal graft. Immediately after surgery, vision will likely be blurry due to corneal swelling, sutures, or the presence of an air/gas bubble (in endothelial keratoplasty). Pain is typically mild to moderate and well-controlled with oral analgesics, gradually resolving over the first few days to a week. As healing progresses and initial swelling subsides, usually within weeks for endothelial grafts and months for full-thickness grafts, vision should steadily improve. The donor cornea should remain transparent, free from haze, edema, or signs of inflammation. Intraocular pressure should remain within a normal range, either spontaneously or with the aid of glaucoma medications if needed. The ocular surface should re-epithelialize completely, and the surgical wound should heal securely without leakage. Patients should experience a resolution of their preoperative symptoms, such as severe pain, chronic irritation, or debilitating photophobia, leading to a significant enhancement in their quality of life and functional vision.

\section{Postoperative Care \& Follow-up}
Postoperative care for corneal transplant patients is rigorous, extensive, and often lifelong, requiring diligent adherence to medication regimens and a structured follow-up schedule to monitor graft health and manage potential complications.

\begin{itemize}
  \item **Frequent Slit-Lamp Examinations:** Initially, visits are daily or weekly, gradually tapering to monthly, quarterly, and eventually annual visits for many years. These examinations meticulously assess graft clarity, signs of inflammation (e.g., keratic precipitates, anterior chamber cells), wound integrity, and suture status (for PKP/DALK).
  \item **Intraocular Pressure (IOP) Checks:** Performed at every visit to monitor for steroid-induced glaucoma or other causes of elevated IOP, which requires prompt management.
  \item **Topical Medication Adjustments:** Dosages of topical corticosteroids, antibiotics, and glaucoma drops are frequently adjusted based on clinical findings and the eye's response.
  \item **Suture Management:** For PKP and DALK, selective suture removal may commence several months post-op, guided by corneal topography to reduce astigmatism. Complete suture removal can take years.
  \item **Refraction and Corrective Lenses:** Once the eye has stabilized (typically 6-12 months post-op), a new refraction is performed, and glasses or contact lenses are prescribed to optimize residual vision.
  \item **Corneal Topography:** Periodically used to monitor corneal shape and astigmatism, providing guidance for suture removal or further refractive procedures.
  \item **Specular Microscopy:** Performed periodically to monitor the endothelial cell count, especially for endothelial grafts, to assess long-term graft health and predict potential late graft failure.
  \item **Rehabilitation:** For patients with residual vision impairment, low vision aids, occupational therapy, or referral to vision rehabilitation specialists may be recommended.
\end{itemize}
Patients are thoroughly educated to recognize and immediately report warning signs such as sudden decrease in vision, increased redness, ocular pain, or light sensitivity, as these symptoms may indicate graft rejection or infection, both of which require urgent medical attention to salvage the graft.

\section{Epidemiology}
Corneal transplantation stands as one of the most frequently performed and successful forms of solid organ transplantation worldwide. In the United States, approximately 40,000 to 50,000 corneal transplants are performed annually, with similar high volumes reported globally. The incidence of corneal diseases necessitating transplantation varies significantly across different geographical regions and demographic groups, influenced by factors such as access to healthcare, prevalence of infectious diseases, and genetic predispositions. Corneal transplants are performed across all age groups, from infants with congenital corneal opacities to elderly individuals suffering from age-related endothelial dystrophies. While there is no significant sex predilection for the procedure itself, certain underlying conditions, such as keratoconus, may show a slight male predominance, whereas Fuchs' endothelial dystrophy is more common in females.

\paragraph{}
The leading indications for corneal transplantation have evolved over time due to advancements in surgical techniques. Historically, full-thickness scarring and keratoconus were the primary drivers for Penetrating Keratoplasty (PKP). Currently, Fuchs' endothelial dystrophy is the most common indication for endothelial keratoplasty (DSEK/DMEK), accounting for a substantial proportion of all transplants performed. Keratoconus remains a significant indication for full-thickness (PKP) or deep anterior lamellar keratoplasty (DALK), particularly in younger patient populations. Other common indications include bullous keratopathy (pseudophakic or aphakic), corneal scarring from various etiologies (e.g., trauma, herpetic keratitis), and regrafting for failed previous transplants. The direct mortality associated with corneal transplant surgery is exceedingly low. Morbidity is primarily related to complications such as graft rejection, primary graft failure, infection, and glaucoma, which can lead to vision loss or the necessity for repeat surgical intervention.

\section*{Outcomes \& Prognosis}
The outcomes of corneal transplantation are generally excellent, with high success rates and significant visual improvement achieved for the vast majority of patients. One-year graft survival rates for Penetrating Keratoplasty (PKP) typically range from 80-90\%, while for endothelial keratoplasty techniques (DSEK/DMEK), they often exceed 90-95\%. Five-year survival rates are lower but remain favorable, generally ranging from 60-80\% for PKP and 80-90\% for endothelial keratoplasty, depending on the underlying disease and specific risk factors. Functional outcomes vary; while many patients achieve good functional vision, often 20/40 or better, achieving 20/20 vision is not universally attainable due to residual astigmatism (especially after PKP), pre-existing ocular comorbidities (e.g., glaucoma, retinal disease), or amblyopia. Endothelial keratoplasty techniques generally offer faster visual recovery, less induced astigmatism, and a lower rejection rate compared to PKP, making them the preferred choice for endothelial dysfunction.

\paragraph{}
Recurrence of the original disease in the donor tissue is rare for most indications, though conditions like herpetic keratitis or certain corneal dystrophies may recur. Graft rejection remains a lifelong risk, necessitating long-term immunosuppression with topical corticosteroids. Prognostic factors significantly influencing outcomes include the degree of corneal vascularization (avascular corneas, such as those in Fuchs' dystrophy or keratoconus, have a better prognosis), the presence of active ocular inflammation, the number of previous graft failures, and patient compliance with the rigorous postoperative medication regimen. Overall, corneal transplantation offers a highly effective and durable means of restoring sight and alleviating pain for a wide spectrum of corneal diseases, profoundly improving patients' quality of life.

\section*{References}
\begin{enumerate}
  \item MedlinePlus. Corneal transplant. U.S. National Library of Medicine; 2025.
  \item Yanoff M, Duker JS. Ophthalmology. 6th ed. Elsevier; 2024.
  \item American Academy of Ophthalmology. Corneal Transplant. EyeWiki; 2024.
  \item Brightbill FS, McDonnell PJ, McGhee CNJ, et al. Corneal Surgery: Theory, Technique and Tissue. 4th ed. Mosby; 2020.
\end{enumerate}

\clearpage